\documentclass[11pt]{article}

\usepackage[letterpaper,margin=1in]{geometry}
\usepackage[T1]{fontenc}
\usepackage{amsmath,amssymb,amsthm,mathtools}
\usepackage{newtxtext,newtxmath}
\usepackage{microtype}
\usepackage{booktabs}
\usepackage{enumitem}
\usepackage[hidelinks]{hyperref}
\usepackage[nameinlink,capitalise,noabbrev]{cleveref}
\usepackage{url}

\allowdisplaybreaks
\numberwithin{equation}{section}
\setlist[itemize]{leftmargin=2em,itemsep=0.15em,topsep=0.25em}

\newtheorem{theorem}{Theorem}[section]
\newtheorem{proposition}[theorem]{Proposition}
\newtheorem{lemma}[theorem]{Lemma}
\newtheorem{corollary}[theorem]{Corollary}
\theoremstyle{remark}
\newtheorem{remark}[theorem]{Remark}

\newcommand{\R}{\mathbb{R}}
\newcommand{\Sph}{\mathbb{S}^{2}}
\newcommand{\Vol}{\operatorname{Vol}}
\newcommand{\Area}{\operatorname{Area}}
\newcommand{\dd}{\,\mathrm d}
\newcommand{\VM}{V_{\mathrm M}}
\newcommand{\arsinh}{\operatorname{arsinh}}
\newcommand{\Gauss}{\operatorname{Gauss}}

\title{Meissner Minimality in the Confocal Peabody Family}
\author{Bruce Swihart\thanks{DCG major-revision checkpoint.  The official-header MPFR replay and independent geometric proofreading remain pre-submission gates.}}
\date{October 8, 2026}

\begin{document}
\maketitle

\begin{abstract}
Arelio, Montejano, and Oliveros introduced a three-parameter family of bodies of constant width obtained from a regular Reuleaux tetrahedron by replacing its six edge neighborhoods with envelopes of spheres associated with confocal quadrics.  The family contains the two classical Meissner bodies as a circle-line degeneration and Robert's tetrahedrally symmetric body as a parabolic limit.  We prove that the Meissner degenerations minimize volume within this regular-tetrahedron confocal family.  We first verify that the explicit eccentricity chart satisfies the source construction for every parameter and every orientation, and we prove the almost-everywhere partition of the normal sphere used in the area bookkeeping.  This gives the exact decomposition
\[
  \Vol K(\mathbf e,\boldsymbol\sigma)
  =\VM+\Phi(e_1)+\Phi(e_2)+\Phi(e_3),
\]
independently of the eight orientation choices.  A complete one-pair calculation reduces \(\Phi\) to a fixed integral with a globally controlled principal-arctangent branch.  The stable chart extends real analytically to the parabolic endpoint.  An independent directed-rounding MPFR certificate proves the high-headroom inequalities
\[
  \Phi(1)>\frac1{4000},
  \qquad
  \Phi''(e)<-\frac1{30000}\quad(0<e<1).
\]
Concavity then gives \(\Phi(e)>e/4000\).  Equality holds exactly when all three parameters vanish.  The result is restricted to regular-tetrahedron confocal Peabodies and does not resolve the three-dimensional Blaschke--Lebesgue problem.
\end{abstract}

\noindent\textbf{Keywords:} bodies of constant width, Peabodies, Meissner bodies, confocal quadrics, strict concavity, MPFR, computer-assisted proof

\section{Introduction}

The Blaschke--Lebesgue problem asks for convex bodies of minimum volume among all convex bodies of prescribed constant width.  In the plane, the minimizers are the Reuleaux triangles \cite{Blaschke1915,Lebesgue1914}.  The corresponding three-dimensional problem remains open.  Bonnesen and Fenchel conjectured that the two noncongruent Meissner bodies are minimizers; their common width-one volume is
\begin{equation}
  \VM
  =\pi\left(\frac23-\frac{\sqrt3}{4}\arccos\frac13\right)
  =0.419860045965080\ldots .
  \label{eq:meissner-volume}
\end{equation}
See \cite{KawohlWeber2011,MartiniMontejanoOliveros2019,Nishioka2026} for background.

Arelio, Montejano, and Oliveros introduced \emph{Peabodies}, obtained by replacing neighborhoods of the six edges of a Reuleaux tetrahedron with sections of envelopes of spheres guided by confocal quadrics \cite{ArelioMontejanoOliveros2023}.  Their regular-tetrahedron construction gives genuine bodies of constant width and contains both Meissner bodies and Robert's tetrahedrally symmetric body.  We study only this regular-tetrahedron family.

The regular tetrahedron has three pairs of opposite edges.  Each pair has a parameter \(e_i\in[0,1)\), and an orientation bit \(\sigma_i\in\{0,1\}\) records which edge receives the elliptic device.  Let
\[
  \mathbf e=(e_1,e_2,e_3),\qquad
  \boldsymbol\sigma=(\sigma_1,\sigma_2,\sigma_3),
\]
and write \(K_d(\mathbf e,\boldsymbol\sigma)\) for the resulting body of width \(d\).  The value \(e_i=0\) is the circle-line Meissner degeneration; the limit \(e_i\to1^-\) is parabolic.

\begin{theorem}[Main theorem]
\label{thm:main}
For every \(d>0\), every \(\mathbf e\in[0,1)^3\), and every \(\boldsymbol\sigma\in\{0,1\}^3\),
\begin{equation}
  \Vol K_d(\mathbf e,\boldsymbol\sigma)
  \ge
  \pi\left(\frac23-\frac{\sqrt3}{4}\arccos\frac13\right)d^3.
  \label{eq:main}
\end{equation}
Equality holds if and only if \(e_1=e_2=e_3=0\).  Over the eight orientation choices, the equality cases form exactly the two classical Meissner congruence classes.
\end{theorem}

The proof has two parts.  First, the source geometry and a partition of the unit normal sphere give an exact additive volume formula.  Second, a one-pair calculation and a compact directed-rounding certificate prove a positive parabolic endpoint and uniform strict concavity.  Positivity then follows from the chord inequality.

\section{Full-parameter Peabody geometry and exact additivity}
\label{sec:additivity}

We work first at the native scale of the source construction: a regular tetrahedron of side length two and a body of constant width two.

\subsection{The explicit chart lies in the source family}

For one opposite-edge pair, let \(0\le e<1\) and put
\begin{equation}
  b^2=\frac{3+3e^2+4\sqrt2\,e}{1-e^2},
  \qquad
  a=\frac b{\sqrt{1-e^2}}.
  \label{eq:ab}
\end{equation}
Since \(b^2\ge3\), define
\[
  \theta=\arccos\frac1b,
  \qquad
  \eta=\arsinh\frac1b,
  \qquad
  u_0=\sin\theta,
  \qquad
  v_0=\cosh\eta.
\]
In an orthonormal frame \((\mathbf k,\mathbf p,\mathbf q)\), set
\begin{align}
  X(t)&=a\sin t\,\mathbf k+b\cos t\,\mathbf p,
  &&\theta\le t\le\pi-\theta,\label{eq:X}\\
  Y(\xi)&=ae\cosh\xi\,\mathbf k+b\sinh\xi\,\mathbf q,
  &&-\eta\le\xi\le\eta.\label{eq:Y}
\end{align}
The curves lie in orthogonal planes, have a common axis, and are confocal because
\[
  a^2-b^2=a^2e^2.
\]
Their transverse endpoint coordinates satisfy
\[
  b\cos\theta=1,
  \qquad
  b\sinh\eta=1,
\]
so the corresponding longitudinal beams are the prescribed opposite edges of length two.

Define the pea radii
\begin{equation}
  R_E(t)=ae(\sin t-u_0),
  \qquad
  R_H(\xi)=a(v_0-\cosh\xi).
  \label{eq:radii}
\end{equation}
They are nonnegative on the parameter rectangle and vanish at the beam endpoints.

\begin{lemma}[Full-parameter admissibility]
\label{lem:admissible}
For every \(e\in(0,1)\), the curves \eqref{eq:X}--\eqref{eq:Y} with radii \eqref{eq:radii} form a pair of convex confocal pea-pod devices satisfying the hypotheses of \cite[Theorem~4.5]{ArelioMontejanoOliveros2023}.  At \(e=0\) the pair is the circle-line degeneration of \cite[Section~5.3]{ArelioMontejanoOliveros2023}.  Interchanging the two devices preserves admissibility.
\end{lemma}

\begin{proof}
Let \(B=b^2\).  Directly from \eqref{eq:ab},
\[
  B+1=\frac{2(e+\sqrt2)^2}{1-e^2},
  \qquad
  B-1=\frac{2(\sqrt2e+1)^2}{1-e^2}.
\]
All factors are positive, so
\[
  \sqrt{B+1}-e\sqrt{B-1}=2\sqrt{1-e^2}.
\]
Since \(v_0=\sqrt{B+1}/b\) and \(u_0=\sqrt{B-1}/b\), this is
\begin{equation}
  a(v_0-eu_0)=2.
  \label{eq:beam-normalization}
\end{equation}
A calculation given in \Cref{sec:pair} shows
\begin{equation}
  |X(t)-Y(\xi)|=a(\cosh\xi-e\sin t).
  \label{eq:center-distance}
\end{equation}
Combining \eqref{eq:radii}, \eqref{eq:beam-normalization}, and \eqref{eq:center-distance} gives
\begin{equation}
  |X-Y|+R_E+R_H=2.
  \label{eq:local-width}
\end{equation}
Thus the explicit chart is precisely a convex confocal device pair in the sense of \cite[Definition~3.5 and Theorem~3.3]{ArelioMontejanoOliveros2023}.  The endpoint statement is \cite[Section~5.3]{ArelioMontejanoOliveros2023}.  Swapping the two devices changes only their labels.
\end{proof}

Apply \Cref{lem:admissible} independently to all three opposite-edge pairs.  An orientation bit merely exchanges the two devices in one pair.  Hence \cite[Theorem~4.5]{ArelioMontejanoOliveros2023} applies for every \(\mathbf e\in[0,1)^3\) and every \(\boldsymbol\sigma\in\{0,1\}^3\): the six wedge-pod patches and four spherical caps bound a convex body of constant width two.

\subsection{Patch decomposition and normal-sphere partition}

For generic parameters, denote the open spherical-cap interiors by
\[
  C_A^\circ,C_B^\circ,C_C^\circ,C_D^\circ
\]
and the six open wedge-pod interiors by
\[
  W_1^{+,\circ},W_1^{-,\circ},\ldots,W_3^{+,\circ},W_3^{-,\circ}.
\]
The complement of these ten open patches is a finite union of seam curves and the four vertices.  The source construction is smooth along nonvertex seams.

For a strictly convex body \(K\), write \(R_K(n)\) for the unique support point with outward normal \(n\in\Sph\).

\begin{lemma}[Opposite support points]
\label{lem:opposite}
For a body of constant width two,
\begin{equation}
  R_K(-n)=R_K(n)-2n.
  \label{eq:opposite-support}
\end{equation}
\end{lemma}

\begin{proof}
Put \(x=R_K(n)\) and \(y=R_K(-n)\).  The supporting planes are separated by two, so \((x-y)\cdot n=2\).  The diameter is two, hence \(|x-y|\le2\).  Equality in Cauchy--Schwarz gives \(x-y=2n\).
\end{proof}

\begin{lemma}[Normal-sphere partition]
\label{lem:normal-partition}
For every generic regular-tetrahedron confocal Peabody, the Gauss images of the ten smooth patch interiors together with the four vertex normal cones partition \(\Sph\) up to spherical measure zero.  Their interiors are pairwise disjoint.  The Gauss images of the seams have spherical area zero.
\end{lemma}

\begin{proof}
A body of constant width is strictly convex.  Therefore each \(n\in\Sph\) has a unique support point.  Two distinct smooth patch interiors cannot have a common outward normal, since the corresponding support plane would touch the body at two points.  Every support point lies in a patch interior, a seam, or a vertex, so the listed sets cover \(\Sph\).

Each seam is a compact piecewise-smooth curve.  Adjacent patches have a common tangent plane along it, so its common normal image is a finite union of curves in \(\Sph\).  Such a set has two-dimensional spherical measure zero.
\end{proof}

Let \(\Omega_A\) be the Gauss image of the radius-two cap centered at \(A\), and let \(N_K(A)\) be the vertex normal cone.

\begin{lemma}[Cap--cone duality]
\label{lem:cap-cone}
For each tetrahedral vertex,
\begin{equation}
  N_K(A)=-\Omega_A.
  \label{eq:cap-cone}
\end{equation}
\end{lemma}

\begin{proof}
If \(x\in C_A^\circ\), then its outward normal is \(n=(x-A)/2\).  By \Cref{lem:opposite}, \(R_K(-n)=A\), so \(-\Omega_A\subset N_K(A)\).  Conversely, if \(m\in N_K(A)\), then \(R_K(-m)=A-2m\) lies on the opposite radius-two cap centered at \(A\), which gives the reverse inclusion.
\end{proof}

For the \(i\)-th opposite-edge pair, let \(\Gamma_i=\Gauss(W_i^+)\).  The binormal pairing of \cite[Lemma~3.8]{ArelioMontejanoOliveros2023} gives
\[
  \Gauss(W_i^-)=-\Gamma_i.
\]
Using \Cref{lem:normal-partition,lem:cap-cone},
\begin{equation}
  4\pi
  =2\sum_{v\in\{A,B,C,D\}}|\Omega_v|
   +2\sum_{i=1}^3|\Gamma_i|.
  \label{eq:normal-partition}
\end{equation}
A patch on a sphere of radius two has area four times the area of its Gauss image.  Therefore the total cap area is
\begin{equation}
  \sum_v|C_v|=8\pi-4\sum_{i=1}^3|\Gamma_i|.
  \label{eq:cap-elimination}
\end{equation}
Adding the wedge-pod areas gives
\begin{equation}
  S_2=8\pi+\sum_{i=1}^3\Psi(e_i),
  \qquad
  \Psi(e_i):=|W_i^+|+|W_i^-|-4|\Gamma_i|.
  \label{eq:additive-area}
\end{equation}
Interchanging the two devices exchanges \(W_i^+\) and \(W_i^-\), so \(\Psi\) is orientation independent.

The generic identity extends to mixed tuples with one or more \(e_i=0\) by the continuous circle-line degeneration in \cite[Section~5.3]{ArelioMontejanoOliveros2023}; equivalently, a collapsed patch may be retained as a zero-area term.

For a body of constant width \(d\), Blaschke's identity is
\begin{equation}
  V=\frac d2S-\frac\pi3d^3.
  \label{eq:blaschke}
\end{equation}
At width two, \eqref{eq:additive-area} gives
\[
  V_2=\frac{16\pi}{3}+\sum_{i=1}^3\Psi(e_i).
\]
Scaling by one half multiplies volume by \(1/8\).  Define
\begin{equation}
  \Phi(e):=\frac{\Psi(e)-\Psi(0)}8.
  \label{eq:phi-def}
\end{equation}
Since the all-zero body is Meissner, we obtain:

\begin{proposition}[Exact additive reduction]
\label{prop:additive}
For every width-one regular-tetrahedron confocal Peabody,
\begin{equation}
  \Vol K_1(\mathbf e,\boldsymbol\sigma)
  =\VM+\Phi(e_1)+\Phi(e_2)+\Phi(e_3),
  \label{eq:additive-volume}
\end{equation}
independently of \(\boldsymbol\sigma\).
\end{proposition}

\subsection{The eight zero-parameter orientations}

\begin{lemma}[Two equality orbits]
\label{lem:orientation-orbits}
At \(\mathbf e=0\), the eight orientation choices form exactly two tetrahedral congruence classes: the four triples of edges incident to a vertex and the four triples bounding a face.  These are the two classical Meissner types.
\end{lemma}

\begin{proof}
Selecting one edge from each of
\[
  \{AB,CD\},\qquad \{AC,BD\},\qquad \{AD,BC\}
\]
gives either a vertex star or a face boundary.  There are four of each.  The tetrahedral symmetry group acts transitively on vertices and on faces, so there are exactly two orbits.
\end{proof}

\section{The one-pair formula}
\label{sec:pair}

We now derive the formula used by the certificate.  Put
\[
  D_0=\cosh\xi-e\sin t,
  \qquad s=\sqrt{1-e^2}.
\]
From \eqref{eq:X}--\eqref{eq:Y},
\begin{align*}
  \frac{|X-Y|^2}{a^2}
  &=(\sin t-e\cosh\xi)^2
    +(1-e^2)(\cos^2t+\sinh^2\xi)\\
  &=\cosh^2\xi-2e\sin t\cosh\xi+e^2\sin^2t\\
  &=(\cosh\xi-e\sin t)^2.
\end{align*}
Since \(D_0>0\) on the parameter rectangle,
\begin{equation}
  d:=|X-Y|=aD_0.
  \label{eq:d}
\end{equation}
Together with \eqref{eq:radii} and \eqref{eq:beam-normalization}, this proves \eqref{eq:local-width}.

Let
\[
  n=\frac{X-Y}{d}=\frac{U}{D_0},
\]
where
\[
  U=(\sin t-e\cosh\xi)\mathbf k
   +s\cos t\,\mathbf p-s\sinh\xi\,\mathbf q.
\]
Since \(|U|=D_0\),
\[
  |n_r|^2=\frac{|U_r|^2-D_{0,r}^2}{D_0^2}.
\]
For \(r=t\),
\[
  U_t=\cos t\,\mathbf k-s\sin t\,\mathbf p,
  \qquad D_{0,t}=-e\cos t,
\]
so \(|U_t|^2-D_{0,t}^2=s^2\).  For \(r=\xi\),
\[
  U_\xi=-e\sinh\xi\,\mathbf k-s\cosh\xi\,\mathbf q,
  \qquad D_{0,\xi}=\sinh\xi,
\]
and again \(|U_\xi|^2-D_{0,\xi}^2=s^2\).  Finally,
\[
  U_t\cdot U_\xi=D_{0,t}D_{0,\xi},
\]
so
\begin{equation}
  n_t\cdot n_\xi=0,
  \qquad
  |n_t|=|n_\xi|=\frac{s}{D_0}=\frac bd.
  \label{eq:conformal}
\end{equation}
Thus
\begin{equation}
  \dd\omega=\frac{b^2}{d^2}\,\dd t\,\dd\xi.
  \label{eq:area-element}
\end{equation}

The paired boundary points are
\[
  U_E=X+R_En,
  \qquad
  U_H=Y-R_Hn.
\]
Using \(X=Y+dn\) and \(d+R_E+R_H=2\),
\[
  U_E=Y+(2-R_H)n,
  \qquad
  U_H=X-(2-R_E)n.
\]
Because \(Y,R_H\) are independent of \(t\), while \(X,R_E\) are independent of \(\xi\),
\begin{align*}
  (U_E)_t&=(2-R_H)n_t,& (U_E)_\xi&=R_En_\xi,\\
  (U_H)_t&=-R_Hn_t,& (U_H)_\xi&=-(2-R_E)n_\xi.
\end{align*}
Consequently,
\[
  \dd A_E=R_E(2-R_H)\frac{b^2}{d^2}\,\dd t\,\dd\xi,
  \qquad
  \dd A_H=R_H(2-R_E)\frac{b^2}{d^2}\,\dd t\,\dd\xi.
\]
Since \(R_E+R_H=2-d\),
\[
  R_E(2-R_H)+R_H(2-R_E)-4=-2(d+R_ER_H).
\]
Therefore
\begin{equation}
  \Psi(e)
  =-2b^2\int_\theta^{\pi-\theta}\int_{-\eta}^{\eta}
  \frac{d+R_ER_H}{d^2}\,\dd\xi\,\dd t.
  \label{eq:psi-2d}
\end{equation}

For fixed \(t\), put \(C=e\sin t\) and \(y=\tanh(\eta/2)\).  The substitution \(z=\tanh(\xi/2)\) gives
\[
  \cosh\xi=\frac{1+z^2}{1-z^2},
  \qquad
  \dd\xi=\frac{2\,\dd z}{1-z^2},
\]
and hence
\begin{equation}
  I(C):=\int_{-\eta}^{\eta}\frac{\dd\xi}{\cosh\xi-C}
  =\frac4{\sqrt{1-C^2}}
  \arctan\left(y\sqrt{\frac{1+C}{1-C}}\right).
  \label{eq:I}
\end{equation}
Differentiating this expression and using
\[
  \frac{2y}{1-y^2}=\sinh\eta=\frac1b
\]
gives
\begin{equation}
  \int_{-\eta}^{\eta}
  \frac{v_0-\cosh\xi}{(\cosh\xi-C)^2}\,\dd\xi
  =\frac{(v_0C-1)I(C)+2/b}{1-C^2}.
  \label{eq:xi-second-integral}
\end{equation}
Substitution into \eqref{eq:psi-2d} and the symmetry \(\sin(\pi-t)=\sin t\) yield
\begin{equation}
\begin{aligned}
  \Psi(e)=-4b^2\int_\theta^{\pi/2}
  \bigg[&\frac{I(e\sin t)}a\\
  &+\frac{e(\sin t-u_0)}{1-e^2\sin^2t}
  \left((v_0e\sin t-1)I(e\sin t)+\frac2b\right)
  \bigg]\dd t.
\end{aligned}
\label{eq:psi-1d}
\end{equation}
At \(e=0\),
\begin{equation}
  \Psi(0)=-\frac{2\pi}{\sqrt3}\arccos\frac13.
  \label{eq:psi-zero}
\end{equation}
Finally set \(x=b\cos t\).  Since \(b\cos\theta=1\), this sends \([\theta,\pi/2]\) to \([1,0]\), and
\[
  A(e,x)=\sin t=\sqrt{1-\frac{x^2}{b^2}},
  \qquad
  \dd t=-\frac{\dd x}{bA}.
\]
Thus
\begin{equation}
  \Psi(e)=\int_0^1H(e,x)\,\dd x,
  \label{eq:fixed-H}
\end{equation}
where
\begin{equation}
H(e,x)=-\frac{4b}{A}
\left[
\frac{I(eA)}a
+\frac{e(A-u_0)}{1-e^2A^2}
\left((v_0eA-1)I(eA)+\frac2b\right)
\right].
\label{eq:fixed-H-explicit}
\end{equation}
The stable parameter
\begin{equation}
  q=\frac{1-e}{1+e},
  \qquad e=\frac{1-q}{1+q}
  \label{eq:q}
\end{equation}
fixes the parabolic endpoint at \(q=0\).  The complete stable formula, branch control, and endpoint analyticity are proved in \Cref{app:q-chart}.

\section{Endpoint value, strict concavity, and scalar positivity}
\label{sec:certificate}

The stable formula extends real analytically to \(0\le q\le1\), and hence \(\Phi\) extends to \([0,1]\).  At \(e=0\), exact differentiation gives
\begin{equation}
\begin{aligned}
  \Phi'(0)=\frac12\bigg[&\frac{5\pi}{3\sqrt3}-3\\
  &+\arccos\frac13
  \left(\frac{3\sqrt2}{2}-\frac{5\sqrt6\,\pi}{18}\right)
  \bigg]>0.
\end{aligned}
\label{eq:phi-prime-zero}
\end{equation}
The endpoint derivative is a consistency check; the compressed proof uses the parabolic endpoint and concavity.

\begin{lemma}[Parabolic endpoint]
\label{lem:parabolic-endpoint}
The endpoint extension satisfies
\begin{equation}
  \Phi(0)=0,
  \qquad
  \Phi(1)>\frac1{4000}.
  \label{eq:endpoint-bound}
\end{equation}
\end{lemma}

\begin{proof}
At \(q=0\), put \(K=(1+\sqrt2)^2\) and \(A=2K+x^2\).  The stable integrand simplifies to
\begin{equation}
\begin{aligned}
H(0,x)=\bigg[&-\frac{16\sqrt2(1+\sqrt2)}{\sqrt A}
 +\frac{4(1-x^2)(A-1)}{A^{3/2}}\bigg]
 \arctan\frac1{\sqrt A}
 -\frac{4(1-x^2)}A.
\end{aligned}
\label{eq:endpoint-integrand}
\end{equation}
The directed-rounding MPFR enclosure described in \Cref{app:mpfr} gives
\[
  \Phi(1)>
  0.0003076528682490342053580258083025\ldots
  >\frac1{4000}.
\]
\end{proof}

\begin{lemma}[Strict concavity]
\label{lem:concavity}
For every \(0<e<1\),
\begin{equation}
  \Phi''(e)<-\frac1{30000}.
  \label{eq:concavity-bound}
\end{equation}
\end{lemma}

\begin{proof}
Let \(\Psi(q)=\int_0^1H(q,x)\,\dd x\).  Since
\[
  q'(e)=-\frac{(1+q)^2}{2},
  \qquad
  q''(e)=\frac{(1+q)^3}{2},
\]
the chain rule gives
\begin{equation}
  \Phi''(e)=\frac{(1+q)^3}{32}
  \int_0^1\left((1+q)H_{qq}(q,x)+2H_q(q,x)\right)\,\dd x.
  \label{eq:concavity-identity}
\end{equation}
The direct MPFR certificate of \Cref{app:mpfr} partitions \([0,1]\) into 32 exact rational slabs and uses ten rational \(x\)-panels on each slab.  Its weakest upper endpoint is
\[
  -0.0000416491903361455479929355233265\ldots
  <-\frac1{30000}.
\]
\end{proof}

\begin{lemma}[Chord bound]
\label{lem:chord}
For every \(0<e<1\),
\begin{equation}
  \Phi(e)>\frac e{4000}>0.
  \label{eq:chord-bound}
\end{equation}
\end{lemma}

\begin{proof}
By \Cref{lem:concavity}, \(\Phi\) is strictly concave.  Therefore it lies above the chord joining its endpoint values:
\[
  \Phi(e)\ge(1-e)\Phi(0)+e\Phi(1).
\]
Apply \Cref{lem:parabolic-endpoint}.
\end{proof}

\section{Proof of the main theorem}

\begin{proof}[Proof of \Cref{thm:main}]
By dilation it is enough to take \(d=1\).  From \Cref{prop:additive,lem:chord},
\[
  \Vol K_1(\mathbf e,\boldsymbol\sigma)-\VM
  =\Phi(e_1)+\Phi(e_2)+\Phi(e_3)\ge0.
\]
Equality forces \(e_1=e_2=e_3=0\), and the converse is immediate.  The two equality congruence classes are identified by \Cref{lem:orientation-orbits}.  Scaling by \(d\) proves \eqref{eq:main}.
\end{proof}

\section{Scope and reproducibility}

The theorem concerns only regular-tetrahedron confocal Peabodies.  It does not prove the global Meissner conjecture, improve a universal lower bound, or establish a six-patch construction for arbitrary semi-regular tetrahedral bases.

The exact part of the proof consists of the source construction, \Cref{lem:admissible,lem:normal-partition,lem:cap-cone,lem:orientation-orbits}, the width normalization, and the one-pair formulas.  The finite proof authority is an independently written C implementation using MPFR directed rounding.  Lower endpoints are evaluated with \texttt{MPFR\_RNDD}, upper endpoints with \texttt{MPFR\_RNDU}; this includes direct calls to \texttt{mpfr\_sqrt} and \texttt{mpfr\_atan}.  The canonical proof uses four endpoint panels and 32 parameter slabs with ten integration panels each.  It is replayed at 384 and 512 bits and in reverse slab order.  A 16-slab under-resolved control and a correction-sector sign mutation are rejected.  The earlier \texttt{mpmath.iv}, Mathematica, and R calculations are retained only as redundant audits.

The public major-revision checkpoint includes the C source, exact rational partitions, all output endpoints, formula-identity checks, proof notes, and replay scripts.  Before submission, the official-header MPFR replay and independent geometric/formula proofreading are required.

\appendix

\section{Stable chart, branch control, and parabolic analyticity}
\label{app:q-chart}

Put
\[
  \kappa=1+\sqrt2,
  \qquad K=\kappa^2=3+2\sqrt2.
\]
From \eqref{eq:q},
\[
  1-e^2=\frac{4q}{(1+q)^2},
  \qquad
  b^2=\frac{K^2+q^2}{2Kq},
  \qquad
  a^2=\frac{(K^2+q^2)(1+q)^2}{8Kq^2}.
\]
Define
\begin{align*}
  D&=\sqrt{K^2+q^2},\\
  R&=\sqrt{K^2+q^2-2Kqx^2},\\
  T&=\sqrt{2K(K^2+q^2)+K^2(1-q)^2x^2}.
\end{align*}
The apparent difference \((R-K+q)/q\) is rationalized using
\[
  R^2-(K-q)^2=2Kq(1-x^2),
\]
which gives
\[
  \delta=\frac{2K(1-x^2)}{R+K-q}.
\]
Likewise,
\[
  R-K=\frac{q(q-2Kx^2)}{R+K}
\]
reduces the second removable quotient to
\[
  M=\frac{K(q-2Kx^2)}{R+K}+(1-K)R-K^2-qR-q-q^2.
\]
Set
\begin{align*}
  P_1&=-\frac{16\sqrt2\,\kappa(K^2+q^2)}{RT},\\
  P_2&=-\frac{4K(1-q^2)(K^2+q^2)\delta M}{RT^3},\\
  Q_0&=-\frac{4K(1-q^2)(K^2+q^2)\delta}{RT^2},\\
  Z&=\frac{K((1+q)D+(1-q)R)}{T(K+q+D)}.
\end{align*}
Then
\begin{equation}
  H(q,x)=(P_1+P_2)\arctan Z+Q_0.
  \label{eq:stable-H}
\end{equation}

We now fix the branch.  In the original formula,
\[
  C=e\sin t\in[0,1),
  \qquad y=\tanh(\eta/2)>0,
\]
so the argument in \eqref{eq:I} is positive and represents an angle in \([0,\pi/2)\).  Put
\[
  A_+=(1+q)D+(1-q)R,
  \qquad
  A_-=(1+q)D-(1-q)R.
\]
A direct calculation gives
\[
  A_+A_-=\frac{2q}{K}T^2,
  \qquad
  K+q-D=\frac{2Kq}{K+q+D}.
\]
Using these positive identities gives, without squaring or changing signs,
\[
  y\sqrt{\frac{1+C}{1-C}}
  =\frac{KA_+}{T(K+q+D)}=Z.
\]
Therefore the principal \(\arctan Z\) in \eqref{eq:stable-H} is the same analytic branch as the geometric angle on the entire domain.

Finally, on \(0\le q,x\le1\),
\[
  D^2\ge K^2,
  \qquad
  R^2\ge(K-q)^2\ge(K-1)^2>0,
  \qquad
  T^2\ge2K^3>0.
\]
Also
\[
  R+K-q\ge2(K-1)>0,
  \quad R+K>0,
  \quad K+q+D>0.
\]
Thus every radicand and denominator is uniformly separated from zero; \(Z>0\); and \eqref{eq:stable-H} contains no negative power of \(q\).  Hence \(H\) is real analytic on a neighborhood of the closed square \([0,1]^2\).  Differentiation under the fixed integral is justified, including on slabs meeting \(q=0\).

\section{Direct MPFR verification}
\label{app:mpfr}

For an interval \([a,b]\), the certifier implements elementary endpoint arithmetic with directed MPFR rounding.  For example,
\[
[a,b]+[c,d]
=[\operatorname{RNDD}(a+c),\operatorname{RNDU}(b+d)],
\]
and multiplication takes the minimum of the four downward-rounded endpoint products and the maximum of the four upward-rounded products.  On positive intervals,
\[
  \sqrt{[a,b]}
  =[\operatorname{RNDD}(\sqrt a),\operatorname{RNDU}(\sqrt b)],
\]
while monotonicity gives
\[
  \arctan([a,b])
  =[\operatorname{RNDD}(\arctan a),\operatorname{RNDU}(\arctan b)].
\]
The implementation calls \texttt{mpfr\_sqrt} and \texttt{mpfr\_atan} directly.

For an \(x\)-panel \([a,b]\), with midpoint \(m\) and width \(h\), the certifier uses the rigorous midpoint enclosure
\begin{equation}
  \int_a^bf(x)\,\dd x
  \in h f(m)+\frac{h^3}{24}f''([a,b]).
  \label{eq:midpoint}
\end{equation}
For the endpoint lemma, four equal rational panels enclose \(\int_0^1H(0,x)\,\dd x\).

For concavity, put
\[
  C(q,x)=(1+q)H_{qq}(q,x)+2H_q(q,x).
\]
A nested automatic-differentiation jet encloses the derivatives through third order in \(q\) and second order in \(x\).  The interval \([0,1]\) is partitioned into
\[
  Q_j=\left[\frac j{32},\frac{j+1}{32}\right],
  \qquad j=0,\ldots,31.
\]
On each slab, centered at \(q_j\), the mean-value form
\[
  \int_0^1C(Q_j,x)\,\dd x
  \subseteq
  \int_0^1C(q_j,x)\,\dd x
  +(Q_j-q_j)\int_0^1C_q(Q_j,x)\,\dd x
\]
is combined with ten rational \(x\)-panels and \eqref{eq:midpoint}.  Multiplication by \((1+Q_j)^3/32\) encloses \(\Phi''\).

The 384-bit forward certificate gives
\[
  \Phi(1)>
  0.0003076528682490342053580258083025767\ldots
  >\frac1{4000},
\]
and
\[
  \sup_{0<e<1}\Phi''(e)
  <-0.0000416491903361455479929355233265745\ldots
  <-\frac1{30000}.
\]
The 512-bit and reverse-order replays give the same reported enclosures.  The 16-slab control and the sign mutation fail, as required.  All 32 MPFR slab enclosures overlap the earlier independent \texttt{mpmath.iv} enclosures.

\begin{thebibliography}{99}
\bibitem{ArelioMontejanoOliveros2023}
I.~Arelio, L.~Montejano, and D.~Oliveros.
\newblock Peabodies of constant width.
\newblock \emph{Beitr. Algebra Geom.}, 64:367--385, 2023.

\bibitem{AnciauxGuilfoyle2011}
H.~Anciaux and B.~Guilfoyle.
\newblock On the three-dimensional Blaschke--Lebesgue problem.
\newblock \emph{Proc. Amer. Math. Soc.}, 139(5):1831--1839, 2011.

\bibitem{Blaschke1915}
W.~Blaschke.
\newblock Konvexe Bereiche gegebener konstanter Breite und kleinsten Inhalts.
\newblock \emph{Math. Ann.}, 76:504--513, 1915.

\bibitem{ChakerianGroemer1983}
G.~D. Chakerian and H.~Groemer.
\newblock Convex bodies of constant width.
\newblock In P.~M. Gruber and J.~M. Wills, editors, \emph{Convexity and Its Applications}, pages 49--96. Birkh\"auser, Basel, 1983.

\bibitem{KawohlWeber2011}
B.~Kawohl and C.~Weber.
\newblock Meissner's mysterious bodies.
\newblock \emph{Math. Intelligencer}, 33(3):94--101, 2011.

\bibitem{Lebesgue1914}
H.~Lebesgue.
\newblock Sur le probl\`eme des isop\'erim\`etres et sur les domaines de largeur constante.
\newblock \emph{Bull. Soc. Math. France}, 42:72--76, 1914.

\bibitem{MartiniMontejanoOliveros2019}
H.~Martini, L.~Montejano, and D.~Oliveros.
\newblock \emph{Bodies of Constant Width: An Introduction to Convex Geometry with Applications}.
\newblock Birkh\"auser, Cham, 2019.

\bibitem{Nishioka2026}
A.~Nishioka.
\newblock An improved lower bound for the three-dimensional Blaschke--Lebesgue problem from spectral and dual perspectives.
\newblock arXiv:2606.01754, 2026.
\end{thebibliography}
\end{document}
